\documentclass[10pt,a4paper]{amsart}
\usepackage[margin=19mm]{geometry}
\usepackage{amsmath,amssymb,mathtools}
\usepackage[scr=rsfs,cal=euler]{mathalfa}
\usepackage{xfrac}
\usepackage[unicode,hidelinks,bookmarks=false]{hyperref}
\newcommand{\ie}{i.e.\ }
\newcommand{\cf}{cf.\ }
\newcommand{\resp}{respectively\ }
\newcommand{\Subsection}{\subsection}
\providecommand{\nobreakdash}{\nobreak\hbox{-}}
\setlength{\emergencystretch}{2em}
\multlinegap0pt
\usepackage{bm}
\usepackage{bbm}
\usepackage{dsfont}
\usepackage{xspace}
\usepackage{cancel}
\usepackage{mathtools}
\usepackage{amsthm}
\makeatletter
\@ifundefined{theorem}{\newtheorem{theorem}{Theorem}}{}
\@ifundefined{corollary}{\newtheorem{corollary}[theorem]{Corollary}}{}
\@ifundefined{definition}{\newtheorem{definition}[theorem]{Definition}}{}
\@ifundefined{lemma}{\newtheorem{lemma}[theorem]{Lemma}}{}
\makeatother
\DeclareMathOperator{\weight}{w}
\newcommand{\cL}{\mathcal{L}}
\newcommand{\cS}{\mathcal{S}}
\newcommand{\gri}{\gamma_{\mathrm{ri},k}}
\newcommand{\setK}{[k]}
\title[A Degree Threshold for Independent Domination in Generalized]{A Degree Threshold for Independent Domination in Generalized Prisms}
\author{Hassine Achour}
\date{Source: arXiv:2608.07956v1 (CC BY 4.0)}
\begin{document}
\maketitle

\noindent\emph{Source and licence.} A Degree Threshold for Independent Domination in Generalized Prisms. Version arXiv:2608.07956v1 (CC BY 4.0), distributed under the Creative Commons Attribution 4.0 International licence, as recorded in the arXiv licence metadata for this exact version and shown on the abstract page https://arxiv.org/abs/2608.07956v1. Full rights notice: licenses/arxiv-2608.07956-CC-BY-4.0.txt.\par\smallskip

\noindent\textit{Editorial scope.} Counted unit: the Theorem whose source label is \texttt{thm:representation} in the article (source0.tex lines 382--398, source label \texttt{thm:representation}), delivered together with the article's own complete original proof of it (source0.tex lines 399--426, one \texttt{proof} environment of 28 lines). No mathematics is written, repaired or re-derived: statement and proof are the article's own text line for line. Required context retained uncounted: def:ridf (lines 147--162); thm:main (lines 293--304); lem:Bnonzero (lines 314--316); cor:lowerbound (lines 321--323). Every definition, display and label of those lines is retained unchanged. Omitted from this package and not redistributed: 1--146, 163--292, 305--313, 317--320, 324--381, 427--580. Nothing inside a retained excerpt is omitted or altered: every retained body is delivered exactly as the article's lines are written, so no omission receipt of any kind accompanies this package. Citations: the retained excerpts carry no citation command at all, so no bibliography entry of the article is transcribed and no reference is redistributed. Reference adaptation: none was needed and none was made; every reference inside the delivered unit resolves inside it, so no unresolved reference or citation is produced. Printed slips kept literal: no printed word, symbol, hypothesis or number of the counted statement or of its proof was altered. Layout and class: the article's own class is \texttt{article} with packages including lmodern, fontenc, inputenc, amsmath, amssymb, amsthm, mathtools, geometry, booktabs, hyperref, graphicx; the wrapper is the lane's pinned \texttt{amsart} editorial wrapper at 10pt on A4, which restates the theorem-like environments the retained text needs and adds the article's own macro definitions that the retained text needs (\texttt{\textbackslash cL}, \texttt{\textbackslash cS}, \texttt{\textbackslash gri}, \texttt{\textbackslash setK}, \texttt{\textbackslash weight}), with extra wrapper packages (bm, bbm, dsfont, xspace, cancel, mathtools). The delivered numbering is the unit's own. Assets: no retained span contains a figure environment, a graphics-inclusion command, a PDF-page inclusion, a table or a TikZ picture; no image file is copied into this package, the package declares no asset dependency and no image file is redistributed with it. Licence: version arXiv:2608.07956v1 of this article is distributed under the Creative Commons Attribution 4.0 International licence, as recorded in the arXiv licence metadata for this exact version and shown on the abstract page https://arxiv.org/abs/2608.07956v1. A full rights notice is delivered at licenses/arxiv-2608.07956-CC-BY-4.0.txt. Characters: every retained excerpt is pure ASCII, so the delivered engine, XeLaTeX with its default Latin Modern text font, reports no missing character.
\begin{definition}[$k$-RiDF, per-colour independent]\label{def:ridf}
Let $k\ge 1$ and $G$ be a graph.
A \emph{per-colour independent $k$-rainbow dominating function} ($k$-RiDF) is a function $f\colon V(G)\to 2^{\setK}$ such that:
\begin{enumerate}
\item[(i)] $|f(v)|\le 1$ for every $v\in V(G)$;
\item[(ii)] for every $i\in\setK$, $V_{i}=\{v:f(v)=\{i\}\}$ is independent in $G$;
\item[(iii)] for every $v$ with $f(v)=\emptyset$, $\bigcup_{u\in N(v)} f(u)=\setK$.
\end{enumerate}
We identify $\emptyset$ with $0$ and $\{i\}$ with $i$.
Write $V_{0}=\{v:f(v)=0\}$ and $V_{i}=\{v:f(v)=\{i\}\}$ for $i\in\setK$, so $(V_{0},V_{1},\dots,V_{k})$ partitions $V(G)$.
Condition~(ii) means \emph{each} colour class $V_{i}$ is independent (not necessarily the union $\bigcup_{i\ge1}V_{i}$); condition~(iii) says every $v\in V_{0}$ has at least one neighbour in each $V_{i}$.
The \emph{weight} is $\weight(f)=\sum_{v\in V(G)}|f(v)|=\sum_{i=1}^{k}|V_{i}|$.
\[
\gri(G)=\min\{\weight(f): f\text{ is a $k$-RiDF of }G\}.
\]
\end{definition}

\begin{theorem}[Main theorem -- hardness at the boundary]\label{thm:main}
For every fixed $k\ge 3$, the language
\[
\cL_{k}=\{H\in\cS_{k}:\gri(H)\le |B|\}
\]
is NP-complete, where $H=S(Q)\in\cS_{k}$ has bipartition $(A,B)=(V(Q),E(Q))$, $\Delta(H)=k$, $H$ is $C_{4}$-free $(k,2)$-biregular, and $B=\{v:\deg(v)=2\}$, so $|B|=|\{v:\deg(v)=2\}|$ is computable from $H$.
Moreover,
\[
\gri(S(Q))\ge |E(Q)|\quad\text{and}\quad \gri(S(Q))=|E(Q)|\iff \chi'(Q)=k.
\]
Hence deciding $\gri(S(Q))=|E(Q)|$ (equivalently $\gri(S(Q))\le|E(Q)|$) is NP-complete under restriction $S(Q)\in\cS_{k}$, with membership in NP.
\end{theorem}

\begin{lemma}[$B$-nonzero]\label{lem:Bnonzero}
Let $H=S(Q)$ with $Q$ $k$-regular, $k\ge 3$, and let $f$ be any $k$-RiDF of $H$. Then $f(b)\neq0$ for every $b\in B$.
\end{lemma}

\begin{corollary}\label{cor:lowerbound}
For $H=S(Q)$ as above, any $k$-RiDF has weight at least $|B|=|E(Q)|$, so $\gri(H)\ge|B|$.
\end{corollary}

\begin{theorem}[Representation theorem]\label{thm:representation}
Let $k\ge 3$ and $Q$ be simple $k$-regular, $H=S(Q)$ with bipartition $(A,B)$.
Define
\[
\mu_{k}(Q)=\min\{|X|: \exists\,c\colon E(Q)\to\setK,\ \ell\colon X\to\setK\text{ such that (i) and (ii) hold}\},
\]
where for $X\subseteq V(Q)=A$:
\begin{itemize}
\item[(i)] for every $a\in A\setminus X$, $\{c(e):e\ni a\}=\setK$ (i.e.\ the $k$ incident edges have pairwise distinct colours);
\item[(ii)] for every $x\in X$, $c(e)\neq\ell(x)$ for every $e\ni x$ (i.e.\ $\ell(x)\notin\{c(e):e\ni x\}$; $x$ misses at least its label colour).
\end{itemize}
Then
\[
\gri(H)=|E(Q)|+\mu_{k}(Q).
\]
Consequently, defining $\delta_{k}(Q)=\gri(S(Q))-|E(Q)|$, we have $\delta_{k}(Q)=\mu_{k}(Q)\ge0$ and $\delta_{k}(Q)=0\iff\chi'(Q)=k$. In particular $0\le\delta_{k}(Q)\le|V(Q)|$.
\end{theorem}

\begin{proof}
Let $f$ be any $k$-RiDF of $H=S(Q)$. By Lemma~\ref{lem:Bnonzero}, $f(b)\neq0$ for all $b\in B$, so each $b$ contributes exactly $1$ (recall $|f(b)|\le 1$): $f(b_{e})=\{c(e)\}$ for some $c(e)\in\setK$ defining $c\colon E(Q)\to\setK$.
Let $X=\{a\in A: f(a)\neq0\}$; for $a\in X$, $f(a)=\{\ell(a)\}$ for some $\ell(a)\in\setK$, and $f(a)=0$ for $a\notin X$.
Weight: $\weight(f)=|B|+|X|=|E(Q)|+|X|$.

We analyse the conditions on $(c,X,\ell)$ imposed by $f$ being a $k$-RiDF.

Independence: each $V_{i}=\{b_{e}:c(e)=i\}\cup\{a\in X:\ell(a)=i\}$ must be independent. Since $B$ is independent and $A$ is independent, the only possible edge inside $V_{i}$ is $a$--$b_{e}$ with $a\in X$, $e\ni a$, $c(e)=i=\ell(a)$.
Hence per-colour independence is equivalent to
\[
c(e)\neq\ell(x)\quad\text{for every }x\in X\text{ and }e\ni x. \tag{$\ast$}
\]

Domination: vertices $b\in B$ and $a\in X$ are non-zero, so need no domination.
For $a\in A\setminus X$ we have $f(a)=0$, so Definition~\ref{def:ridf}(iii) requires $\bigcup_{b\in N_{H}(a)}f(b)=\setK$, i.e.
\[
\{c(e): e\ni a\}=\setK. \tag{$\dagger$}
\]
This uses that $a$ has exactly $k$ incident edges: the union has size $k$ (all colours) iff the $k$ singletons are distinct. No condition is imposed on $a\in X$ (already dominated by being non-zero).
Thus every $k$-RiDF yields a triple satisfying ($\dagger$) for $a\notin X$ and ($\ast$) for $x\in X$, with weight $|E|+|X|$.

Conversely, given $(c,X,\ell)$ satisfying ($\dagger$) and ($\ast$), define $f(b_{e})=\{c(e)\}$, $f(a)=0$ for $a\notin X$, $f(a)=\{\ell(a)\}$ for $a\in X$.
Then ($\ast$) ensures each $V_{i}$ is independent, ($\dagger$) ensures every $a\notin X$ sees all $k$ colours, and all other vertices are non-zero. Hence $f$ is a $k$-RiDF with weight $|B|+|X|$.

Therefore $\gri(H)=\min_{f}\weight(f)=|B|+\min|X|=|E(Q)|+\mu_{k}(Q)$ with $\mu_{k}(Q)$ as defined. The minimum is over triples satisfying exactly the two bullet conditions above; the stronger exact-miss condition $\{c(e):e\ni x\}=\setK\setminus\{\ell(x)\}$ is \emph{not} required. A labelled vertex may miss several colours, e.g.\ all incident edges could share one colour different from $\ell(x)$, which still satisfies ($\ast$) and is a valid local configuration.

Defining $\delta_{k}(Q)=\gri(S(Q))-|E(Q)|$ we obtain $\delta_{k}(Q)=\mu_{k}(Q)$. By Corollary~\ref{cor:lowerbound}, $\delta_{k}(Q)\ge0$, and $\delta_{k}(Q)=0\iff\chi'(Q)=k$ by Theorem~\ref{thm:main}.
\end{proof}

\end{document}
